\documentclass{situationist-preprint}

\title{The Pessimistic Meta-Induction Reconsidered}
\msid{SCI-02-b}
\msfamily{SCI-02}
\msdiscipline{Philosophy of Science}
\mstype{Companion manuscript}
\mskeywords{pessimistic meta-induction; scientific realism; approximate truth; base-rate fallacy; theory change}

\begin{document}
\maketitle

\begin{abstract}
SCI-02-a defends scientific realism and answers the pessimistic meta-induction by showing that the falsity of past theories does not entail that they were not approximately true. Here we reconsider the meta-induction in its two main forms, the inductive and the reductio form, and show that each needs a premise its historical evidence does not supply. The inductive form needs a base rate of mature theories with false essential posits; Laudan's list does not provide one, because most of its members were not mature and the false posits of the mature ones were idle. The reductio form needs the premise that past essential posits were false, and the record of retention shows that past theories were partially true.
\end{abstract}

\section{Two forms of the argument}\label{s:forms}

\citet{Laudan1981} offered a list of theories that were successful in their time and are now regarded as false, among them the crystalline spheres of ancient astronomy, the humoral theory of medicine, the effluvial theory of static electricity, the phlogiston theory of chemistry, the caloric theory of heat, and the optical and electromagnetic ethers. Two arguments have been drawn from the list.

The \emph{inductive form}: most successful past theories were false; current theories are relevantly like past ones; so current theories are probably false. The \emph{reductio form}: if realism were true, success would be a reliable sign of approximate truth; past theories were successful and not approximately true; so success is not a reliable sign, and realism, which infers truth from success, is unsupported \citep{Lewis2001}.

SCI-02-a supplies the terms for assessing both. A theory is mature if it has at least one novel success (Definition~3.4); its essential posits are the claims used ineliminably in deriving its novel successes (Definition~3.2); and it is approximately true if its essential posits are true (Definition~3.3). Realism, in the form of Theorem~4.2 of SCI-02-a, is the claim that every mature theory is approximately true.

\section{The inductive form}\label{s:inductive}

\begin{proposition}\label{p:base}
The inductive form supports the conclusion that current mature theories have false essential posits only if the proportion of past mature theories with false essential posits is high. Laudan's list does not show that it is.
\end{proposition}

\begin{proof}
The first claim follows from the logic of statistical inference: an induction to a property of the members of a class requires a sample from that class. The relevant class is mature theories, and the relevant property is having a false essential posit, since falsity elsewhere in a theory is compatible with approximate truth (Proposition~4.3 of SCI-02-a). The list does not sample that class with respect to that property. Most of its entries, the crystalline spheres, humoral medicine and the effluvial theory among them, had no novel successes and so are not mature. For the entries with novel successes, the false posits were not essential: the derivation of Fresnel's novel predictions from his theory used the laws governing the propagation of light, not the mechanical constitution of the ether \citep{Psillos1999}, and the predictive successes of caloric theory were derived from claims about the conservation and flow of heat that survive in thermodynamics.
\end{proof}

The inductive form also commits the error that \citet{MagnusCallender2004} identify in many arguments about realism: it infers a probability for current theories from the frequency of falsity in a historical sample without attending to the base rate of approximately true theories in the population from which current theories are drawn. Since the population of current mature theories differs from the population of past theories, which included many immature ones, the historical frequency is not that base rate.

\section{The reductio form}\label{s:reductio}

The reductio form avoids statistical inference, and \citet{Lewis2001} showed that it requires only that some past successful theories be false, not that most were. But its second premise is not that past successful theories were false; it is that they were not approximately true. Proposition~4.3 of SCI-02-a shows that the first does not entail the second. The premise must be established theory by theory, by showing that an essential posit of a mature past theory was false.

\begin{proposition}\label{p:retention}
For each sequence of mature theories in which the successor retains the essential content of its predecessor in the sense of Definition~3.8 of SCI-02-a, the approximate truth of the latest member implies the approximate truth of every earlier member; and the historical cases on Laudan's list for which retention has been established fall under this condition.
\end{proposition}

\begin{proof}
The first part is Corollary~4.5 of SCI-02-a. The second part is a historical claim: the transitions from Fresnel's theory to Maxwell's and from caloric theory to thermodynamics are the cases for which retention of the posits used in deriving novel predictions has been documented \citep{Psillos1999}, and they are the central mature cases on the list.
\end{proof}

Proposition~\ref{p:retention} turns the history against the meta-induction. The pessimist reads the record of theory change as a record of discarded theories. Read with the distinction between essential and idle posits, it is a record of retained essential content: past theories were false, but partially true, and the part that was true is the part their successes depended on.

\section{Objections and replies}\label{s:objections}

\begin{objection}[Hindsight]
Essential posits are identified with hindsight: whatever survives is declared essential, so retention is guaranteed by the method of identification \citep{Stanford2006}.
\end{objection}

\begin{reply}
Definition~3.2 of SCI-02-a identifies essential posits by the structure of the derivation of a novel success, which is available when the success is derived and does not depend on what later theories say. \citet{Vickers2013} shows how the posits that do the work in a derivation can be identified case by case from the derivation itself. The test can fail: had Fresnel's derivations required a claim about the ether's mechanical constitution that his successors denied, the case would count against realism.
\end{reply}

\begin{objection}[Recalcitrant cases]
There are cases in which false claims appear to have been used essentially, such as Kirchhoff's theory of diffraction \citep{SaatsiVickers2011}.
\end{objection}

\begin{reply}
SCI-02-a answers this objection under the heading of false essential posits. We add only that the reductio form, to succeed, needs such cases to be the rule among mature theories and not the exception, since a single unresolved case does not show that success is generally unreliable as a sign of approximate truth.
\end{reply}

\section{Conclusion}\label{s:conclusion}

In its inductive form the meta-induction samples the wrong class with respect to the wrong property; in its reductio form it relies on a premise the historical record, read with the distinction between essential and idle posits, does not support. Past mature theories were false, but they were partially true, and the meta-induction fails, as SCI-02 concludes.

\begin{thebibliography}{9}
\bibitem[Laudan(1981)]{Laudan1981} Laudan, L. (1981). A confutation of convergent realism. \emph{Philosophy of Science}, 48(1), 19--49.
\bibitem[Lewis(2001)]{Lewis2001} Lewis, P.~J. (2001). Why the pessimistic induction is a fallacy. \emph{Synthese}, 129(3), 371--380.
\bibitem[Magnus and Callender(2004)]{MagnusCallender2004} Magnus, P.~D., and Callender, C. (2004). Realist ennui and the base rate fallacy. \emph{Philosophy of Science}, 71(3), 320--338.
\bibitem[Psillos(1999)]{Psillos1999} Psillos, S. (1999). \emph{Scientific Realism: How Science Tracks Truth}. Routledge.
\bibitem[Saatsi and Vickers(2011)]{SaatsiVickers2011} Saatsi, J., and Vickers, P. (2011). Miraculous success? Inconsistency and untruth in Kirchhoff's diffraction theory. \emph{The British Journal for the Philosophy of Science}, 62(1), 29--46.
\bibitem[Stanford(2006)]{Stanford2006} Stanford, P.~K. (2006). \emph{Exceeding Our Grasp: Science, History, and the Problem of Unconceived Alternatives}. Oxford University Press.
\bibitem[Vickers(2013)]{Vickers2013} Vickers, P. (2013). A confrontation of convergent realism. \emph{Philosophy of Science}, 80(2), 189--211.
\end{thebibliography}

\end{document}
